% ════════════════════════════════════════════
% 清洗规则（5.2.3）
% ════════════════════════════════════════════
\subsubsection{清洗规则}

本节把清洗规则写成可执行的变换，并说明每条规则对后续回归的影响。

缺失填补以完全随机缺失为前提：若缺失只与未被观测的偶然因素有关，缺失帖子与观测帖子来自同一分布，就可以用观测样本的众数作为缺失值的替代。本样本中付费标记的众数为 0，即把缺失帖子视为自然帖子；填补只修改这一列，并附加一个填补标记，便于追溯与敏感性对照。虽然本部分只有 1 条缺失，保留标记仍有助于在结果中区分原始值与填补值。

编码采用“K 个水平保留 K-1 个哑变量”的规则，舍弃的水平作为参照类。帖子类型、内容主题、月份、星期与时段的水平数分别为 4、3、12、7、3，加上常数项与付费标记，最终模型共需 $1+3+2+11+6+2+1=26$ 个参数。若不丢弃参照类，设计矩阵将完全共线，最小二乘解不唯一；因此参照类的选择只影响系数的解释口径，不影响模型拟合。

时段合并是偏差与方差的权衡。逐小时编码保留时间分辨率，但每个小时平均只有约 20 条帖子，最稀疏的单元可能仅有个位数观测，系数标准误会被放大。合并后凌晨、上午、下午晚间三组分别有 196、193、111 条，均超过 100 条，满足稳定估计的基本要求；由此带来的设定偏差，通过后续残差检验与稳健性对比间接评估。

对数变换后，系数不再直接等于曝光次数的边际变化，但在对数尺度上具有近似百分比含义。对连续变量，系数乘以 100 约等于解释变量增加一个单位时曝光变化的百分比；对哑变量，取指数直接给出相对参照类的倍数。当系数绝对值小于 0.2 时，指数减一与系数本身接近，正文同时报告两种口径以免误读。变换的合理性由偏度与峰度的下降幅度衡量：原始曝光量与 log1p 曝光量的偏度之差、峰度之差越大，说明变换对尾部的压缩越明显。

填补方案是否影响结论，用两种方案下付费系数的绝对差判断。若差值小于 0.01，则认为填补不改变参数的经济含义与显著性判断；否则需要在正文中同时报告两种结果并说明取舍。

综上所述，本节完成了缺失填补、独热编码、时段合并与对数变换的规则设计，得到了包含 26 个参数的分析数据框架，接下来将基于此给出清洗结果与敏感性检验。
